\cgeentry{cge_variable}{User variables, defined by an expression or determined by user equations}

\cgesec{Description}

\cmd{cge_variable} adds a variable to the calibrated model. With an expression,
it adds the variable and the equation that defines it; with \opt{}{value()}, it
adds a free variable, with that base value, that user equations added by
\manref{cge_equation} determine. A variable may be indexed by a set of the
model, and is then defined for every element of the set.

\cgesec{Quick start}

Government purchases at market prices
\begin{stcode}
cge_variable GEXP = sum(c, PQ(c) * QG(c)), nominal
\end{stcode}
Fix them in value, with the volumes adjusting
\begin{stcode}
cge_closure, swap(GEXP GADJ)
\end{stcode}
Prices relative to the consumer price index, for every commodity
\begin{stcode}
cge_variable PREL(c) = PQ(c) * (1 + tv(c)) / CPI
\end{stcode}
A variable with base value 1, for an equation of \manref{cge_equation} to
determine
\begin{stcode}
cge_variable LAM, value(1)
\end{stcode}

\cgesec{Syntax}

A variable and its definition
\begin{stsyntax}
\cmdb{cge_variable} \argu{name}[\cmd{(}\argu{i}\cmd{)}] \cmd{=} \argu{expression}
[\texttt{,} \opt{}{nominal}]
\end{stsyntax}
A variable that user equations determine
\begin{stsyntax}
\cmdb{cge_variable} \argu{name}[\cmd{(}\argu{i}\cmd{)}]\texttt{,} \optarg{}{value}{\#}
[\opt{}{nominal}]
\end{stsyntax}
\argu{name} has at most 24 characters, letters, digits and underscores; it
may not be the name of a model variable, a user variable or a parameter that
expressions use, and a defined variable may not take the name of a user
equation. \argu{i} is an index: the letter of a set, or the name of a user set
(\manref{cge_set}), possibly followed by digits, as in \cmd{c} or \cmd{c2}. \argu{expression} is written in the language
described in \manref{cge_equation}.

\begin{opttable}
\optgroup{Main}
\optarg{}{value}{\#} & base value of a variable that user equations determine\\
\opt{}{nominal} & the variable is measured in local currency\\
\end{opttable}

\cgesec{Options}

\begin{optdesc}
\item[\optarg{}{value}{\#}] gives the base value of a variable that user
equations determine, and adds no equation. A variable has either a definition
or a value, not both.

\item[\opt{}{nominal}] marks a variable measured in local currency, so that
the homogeneity check of \manref{cge_check} expects it to scale with the
numeraire.
\end{optdesc}

\cgesec{Remarks and examples}

User variables extend the model without changing its code. A defined variable
gives a name to an expression of the model's variables, which can then be
reported, bounded by \manref{cge_constraint} or, through a swap, fixed in the
closure. A variable given a value is a new unknown, and needs equations of its
own.

User variables are free when they are added. A defined variable is determined
by its definition. Swapped into the closure with
\cmd{cge_closure, swap(}\argu{name} \argu{variable}\cmd{)}, it is fixed, which
fixes the expression that defines it, and \argu{variable}, which the closure
fixed, adjusts instead; \manref{cge_shock} can then change its value. A
variable given a value adds one free variable for each element and no
equation, so the closure is not square until \manref{cge_equation} adds as
many equations; until then \manref{cge_solve} stops with error 459.

Indexes name the sets of the model: \cmd{a} activities, \cmd{c} commodities,
\cmd{f} factors, \cmd{h} households, \cmd{d} households and enterprises,
\cmd{i} institutions that receive factor income and \cmd{n} nests of value
added. \argu{name}\cmd{(}\argu{i}\cmd{)} defines the variable for every element
of the set of \argu{i}, and the index may then appear in the expression.

User variables are saved with the model by \manref{cge_save}, kept when the
model is calibrated again, with base values from the new base, solved in every
simulation, dynamic run and sensitivity draw, listed in the results with the
other variables, and written into the GAMS program by \manref{cge_gams}.
Simulations stored before a variable was added hold its base value until they
are solved again. \cmd{cge_equation, list} lists the user variables, and
\cmd{cge_equation, clear} removes them with the user equations.

\cgesubsec{Example 1: government spending fixed in value}

The synthetic SAM is mapped and calibrated as in \manref{cge_solve}. The value
of government purchases at market prices is defined, swapped for the scaling
factor of government consumption, and so fixed while world import prices rise
by 10 per cent:
\stlog{var_gexp}
The value of purchases stays at its base, 45, and the volumes adjust, all in
the same proportion:
\stlog{var_gexp_res}

\cgesubsec{Example 2: an indexed variable}

A variable indexed by commodity is defined for each one; \cmd{tv(c)}, the VAT
rate, is zero in this SAM. Solving the simulation again gives the variable its
values in the solution:
\stlog{var_index}

\cgesubsec{Example 3: a rigid real wage}

Fixing the real wage of labour, and freeing its supply, gives unemployment at
a fixed real wage. The shock and the swap of Example~1 are removed first:
\stlog{var_rw}
Under the CPI numeraire the real wage and the wage are both fixed at the base,
and the rise in import prices lowers employment by about 5 per cent.

\cgesubsec{Example 4: a variable that equations determine}

A variable given a value has no equation, and the closure is then not square,
as \manref{cge_check} shows; \manref{cge_solve} stops with error 459 before it
solves:
\stlog{var_value}
\manref{cge_equation} gives an example that adds the equation.

\cgesec{Stored results}

\cmd{cge_variable} stores the following in \cmd{r()}:
\begin{storedres}
\resgroup{Scalars}
\res{r(N_elements)} & number of elements of the variable\\
\resgroup{Macros}
\res{r(name)} & name of the variable\\
\end{storedres}

\cgesec{Methods and formulas}

A defined variable $U$ with index $i$ adds, for every element $e$ of the set of
$i$, the variable $U_e$ and the equation
\begin{equation*}
U_e = g_e(x),
\end{equation*}
where $g_e$ is the expression with $i$ replaced by $e$ and $x$ the vector of the
model's variables, parameters at their current values. Its base value is
$U^0_e = g_e(x^0)$, which must not be missing. The equation is a user equation
named after the variable, and \cmd{cge_equation, list} lists it with the
others. A variable given a value $v$ adds $U_e$ with $U^0_e = v$ and no
equation. User variables come after the model's variables in the stacked
vector of variables, in the order in which they are added, so that a
definition may use the variables added before it. \Chapref{chap:user} describes
the compilation of the expressions and their derivatives.

\cgesec{Also see}

\begin{alsosee}
\seeentry{cge_equation}{User equations: new relations among the model's variables}
\seeentry{cge_closure}{Choose the model closure}
\seeentry{cge_shock}{Change a parameter or a fixed variable}
\seeentry{cge_gams}{Write the model as a GAMS program, and solve it in GAMS}
\seechapter{chap:user}{User variables and equations}
\end{alsosee}
